% GENERATED FILE. Edit canonical lessons, then run npm run generate:books.
% canonical-source-hash: fnv1a64:4653d1dc2c21f471
% canonical-lessons: TE-C451-riksa, TE-C451-edlabandi, TE-C451-topudubandi, TE-C451-kalinadaka, TE-C451-batasari

\chapter{Old Ways of Getting About}
\label{ch:te-old-ways-of-getting-about}
\hlchaptermodality{451}

\begin{chapteropening}
\textbf{By the end of this chapter:} \emph{I can say a rickshaw, a bullock cart, a handcart, going on foot and a wayfarer.}

Complete the last lesson of chapter 451: I can say a rickshaw, a bullock cart, a handcart, going on foot and a wayfarer.
\end{chapteropening}

\section[\texorpdfstring{\te{రిక్షా} (rikṣā)}{రిక్షా (rikṣā)}]{\te{రిక్షా} (rikṣā) — a rickshaw}

\label{lesson:TE-C451-riksa}



\begin{quote}
Before the new one: say the Telugu for to unload, then the Telugu for to drop someone off.
\end{quote}

\subsection*{You'll want to know: \te{రిక్షా}}
\textbf{\te{రిక్షా}} — \emph{rikṣā} — \textquotedblleft{}a rickshaw\textquotedblright{}.

\begin{quote}
\te{ఆ} \te{రోజుల్లో} \te{మా} \te{తాత} \te{రిక్షా} \te{తొక్కేవాడు}.
\end{quote}

\begin{quote}
\emph{ā rōjullō mā tāta rikṣā tokkēvāḍu.}
\end{quote}

\begin{quote}
In those days my grandfather used to pedal a rickshaw.
\end{quote}

A cycle rickshaw is pedalled: \te{రిక్షా} \te{తొక్కు}. The motor \te{ఆటో} took its place in most towns.

\begin{cousinweb}
\textbf{\te{రిక్షా}} (\emph{rikṣā}) is a loan from English \emph{rickshaw}, which English shortened from a Japanese word for a man-pulled cart.
\end{cousinweb}

\begin{grammarlens}[title={What you've built}]
One more word for getting around.
\end{grammarlens}

\subsection*{Your turn}
\begin{itemize}
\raggedright
  \item \emph{Say it:} \emph{rikṣā}
  \item \emph{Say it:} \emph{rikṣā}, once more
  \item \emph{Say it:} \emph{ā rōjullō mā tāta rikṣā tokkēvāḍu.} — In those days my grandfather used to pedal a rickshaw.
  \item \emph{Recall:} say the Telugu for to unload, then the Telugu for to drop someone off, then say \emph{rikṣā} again
\end{itemize}

\begin{tcolorbox}[breakable,colback=teal!4,colframe=teal!35!black,title={Before you move on}]
What does \te{రిక్షా} mean? (\textquotedblleft{}A rickshaw\textquotedblright{}.) Say it once more.
\end{tcolorbox}
\section[\texorpdfstring{\te{ఎడ్లబండి} (eḍlabaṇḍi)}{ఎడ్లబండి (eḍlabaṇḍi)}]{\te{ఎడ్లబండి} (eḍlabaṇḍi) — a bullock cart}

\label{lesson:TE-C451-edlabandi}



\begin{quote}
Before the new one: say the Telugu for to drop someone off, then the Telugu for a rickshaw.
\end{quote}

\subsection*{You'll want to know: \te{ఎడ్లబండి}}
\textbf{\te{ఎడ్లబండి}} — \emph{eḍlabaṇḍi} — \textquotedblleft{}a bullock cart\textquotedblright{}.

\begin{quote}
\te{పల్లెల్లో} \te{ఇంకా} \te{ఎడ్లబండ్లు} \te{కనిపిస్తాయి}.
\end{quote}

\begin{quote}
\emph{pallellō iṅkā eḍlabaṇḍlu kanipistāyi.}
\end{quote}

\begin{quote}
Bullock carts are still seen in the villages.
\end{quote}

\te{ఎడ్ల}, of oxen, + \te{బండి}, a cart. Plural \te{ఎడ్లబండ్లు}.

\begin{cousinweb}
\textbf{\te{ఎడ్లబండి}} (\emph{eḍlabaṇḍi}) is a Telugu compound: \te{ఎడ్ల}, from \te{ఎద్దు}, \textquotedblleft{}an ox\textquotedblright{}, plural \te{ఎడ్లు}, native Dravidian (DEDR 815), and \te{బండి}, \textquotedblleft{}a cart\textquotedblright{}, which DEDR lists among its Indo-Aryan loans (DEDR App. 50).
\end{cousinweb}

\begin{grammarlens}[title={What you've built}]
One more word for getting around.
\end{grammarlens}

\subsection*{Your turn}
\begin{itemize}
\raggedright
  \item \emph{Say it:} \emph{eḍlabaṇḍi}
  \item \emph{Say it:} \emph{eḍlabaṇḍi}, once more
  \item \emph{Say it:} \emph{pallellō iṅkā eḍlabaṇḍlu kanipistāyi.} — Bullock carts are still seen in the villages.
  \item \emph{Recall:} say the Telugu for to drop someone off, then the Telugu for a rickshaw, then say \emph{eḍlabaṇḍi} again
\end{itemize}

\begin{tcolorbox}[breakable,colback=teal!4,colframe=teal!35!black,title={Before you move on}]
What does \te{ఎడ్లబండి} mean? (\textquotedblleft{}A bullock cart\textquotedblright{}.) Say it once more.
\end{tcolorbox}
\section[\texorpdfstring{\te{తోపుడుబండి} (tōpuḍubaṇḍi)}{తోపుడుబండి (tōpuḍubaṇḍi)}]{\te{తోపుడుబండి} (tōpuḍubaṇḍi) — a handcart; a pushcart}

\label{lesson:TE-C451-topudubandi}



\begin{quote}
Before the new one: say the Telugu for a rickshaw, then the Telugu for a bullock cart.
\end{quote}

\subsection*{You'll want to know: \te{తోపుడుబండి}}
\textbf{\te{తోపుడుబండి}} — \emph{tōpuḍubaṇḍi} — \textquotedblleft{}a handcart; a pushcart\textquotedblright{}.

\begin{quote}
\te{అతను} \te{తోపుడుబండి} \te{మీద} \te{పండ్లు} \te{అమ్ముతాడు}.
\end{quote}

\begin{quote}
\emph{atanu tōpuḍubaṇḍi mīda paṇḍlu ammutāḍu.}
\end{quote}

\begin{quote}
He sells fruit from a handcart.
\end{quote}

\te{తోపుడు}, pushing, + \te{బండి}, a cart: a cart pushed by hand. Street sellers use them.

\begin{cousinweb}
\textbf{\te{తోపుడుబండి}} (\emph{tōpuḍubaṇḍi}) is a Telugu compound: \te{తోపుడు}, \textquotedblleft{}pushing\textquotedblright{}, from the native verb \emph{trōcu}, \textquotedblleft{}to push, to shove\textquotedblright{} (DEDR 3340), and \te{బండి}, \textquotedblleft{}a cart\textquotedblright{} (DEDR App. 50).
\end{cousinweb}

\begin{grammarlens}[title={What you've built}]
One more word for getting around.
\end{grammarlens}

\subsection*{Your turn}
\begin{itemize}
\raggedright
  \item \emph{Say it:} \emph{tōpuḍubaṇḍi}
  \item \emph{Say it:} \emph{tōpuḍubaṇḍi}, once more
  \item \emph{Say it:} \emph{atanu tōpuḍubaṇḍi mīda paṇḍlu ammutāḍu.} — He sells fruit from a handcart.
  \item \emph{Recall:} say the Telugu for a rickshaw, then the Telugu for a bullock cart, then say \emph{tōpuḍubaṇḍi} again
\end{itemize}

\begin{tcolorbox}[breakable,colback=teal!4,colframe=teal!35!black,title={Before you move on}]
What does \te{తోపుడుబండి} mean? (\textquotedblleft{}A handcart; a pushcart\textquotedblright{}.) Say it once more.
\end{tcolorbox}
\section[\texorpdfstring{\te{కాలినడక} (kālinaḍaka)}{కాలినడక (kālinaḍaka)}]{\te{కాలినడక} (kālinaḍaka) — going on foot}

\label{lesson:TE-C451-kalinadaka}



\begin{quote}
Before the new one: say the Telugu for a bullock cart, then the Telugu for a handcart.
\end{quote}

\subsection*{You'll want to know: \te{కాలినడక}}
\textbf{\te{కాలినడక}} — \emph{kālinaḍaka} — \textquotedblleft{}going on foot\textquotedblright{}.

\begin{quote}
\te{మేము} \te{కొండ} \te{మీదకు} \te{కాలినడకన} \te{వెళ్ళాం}.
\end{quote}

\begin{quote}
\emph{mēmu koṇḍa mīdaku kālinaḍakana veḷḷāṁ.}
\end{quote}

\begin{quote}
We went up the hill on foot.
\end{quote}

\te{కాలినడకన}, on foot, is the usual form. \te{కాలి}, of the foot, + \te{నడక}, walking.

\begin{cousinweb}
\textbf{\te{కాలినడక}} (\emph{kālinaḍaka}) is native Dravidian: \te{కాలి}, from \te{కాలు}, \textquotedblleft{}foot\textquotedblright{} (DEDR 1479), and \te{నడక}, \textquotedblleft{}walking\textquotedblright{}, the noun of \te{నడుచు}, \textquotedblleft{}to walk\textquotedblright{} (DEDR 3582).
\end{cousinweb}

\begin{grammarlens}[title={What you've built}]
One more word for getting around.
\end{grammarlens}

\subsection*{Your turn}
\begin{itemize}
\raggedright
  \item \emph{Say it:} \emph{kālinaḍaka}
  \item \emph{Say it:} \emph{kālinaḍaka}, once more
  \item \emph{Say it:} \emph{mēmu koṇḍa mīdaku kālinaḍakana veḷḷāṁ.} — We went up the hill on foot.
  \item \emph{Recall:} say the Telugu for a bullock cart, then the Telugu for a handcart, then say \emph{kālinaḍaka} again
\end{itemize}

\begin{tcolorbox}[breakable,colback=teal!4,colframe=teal!35!black,title={Before you move on}]
What does \te{కాలినడక} mean? (\textquotedblleft{}Going on foot\textquotedblright{}.) Say it once more.
\end{tcolorbox}
\section[\texorpdfstring{\te{బాటసారి} (bāṭasāri)}{బాటసారి (bāṭasāri)}]{\te{బాటసారి} (bāṭasāri) — a wayfarer}

\label{lesson:TE-C451-batasari}



\begin{quote}
Before the new one: say the Telugu for a handcart, then the Telugu for going on foot.
\end{quote}

\subsection*{You'll want to know: \te{బాటసారి}}
\textbf{\te{బాటసారి}} — \emph{bāṭasāri} — \textquotedblleft{}a wayfarer\textquotedblright{}.

\begin{quote}
\te{ఆ} \te{చెట్టు} \te{నీడలో} \te{బాటసారులు} \te{కూర్చునేవారు}.
\end{quote}

\begin{quote}
\emph{ā ceṭṭu nīḍalō bāṭasārulu kūrcunēvāru.}
\end{quote}

\begin{quote}
Wayfarers used to sit in the shade of that tree.
\end{quote}

Plural: \te{బాటసారులు}, \emph{bāṭasārulu}.

\te{బాట}, a path, + -\te{సారి}, one who goes. An old-fashioned, warm word for a traveller on foot.

\begin{cousinweb}
\textbf{\te{బాటసారి}} (\emph{bāṭasāri}) is a Telugu compound: \te{బాట}, \textquotedblleft{}a path\textquotedblright{} (CDIAL 11366), and -\te{సారి}, \textquotedblleft{}one who goes\textquotedblright{}, usually traced to Sanskrit \emph{-sārin}, \textquotedblleft{}going, moving\textquotedblright{}.
\end{cousinweb}

\begin{grammarlens}[title={What you've built}]
One more word for getting around.
\end{grammarlens}

\subsection*{Your turn}
\begin{itemize}
\raggedright
  \item \emph{Say it:} \emph{bāṭasāri}
  \item \emph{Say it:} \emph{bāṭasāri}, once more
  \item \emph{Say it:} \emph{ā ceṭṭu nīḍalō bāṭasārulu kūrcunēvāru.} — Wayfarers used to sit in the shade of that tree.
  \item \emph{Recall:} say the Telugu for a handcart, then the Telugu for going on foot, then say \emph{bāṭasāri} again
\end{itemize}

\begin{tcolorbox}[breakable,colback=teal!4,colframe=teal!35!black,title={Before you move on}]
What does \te{బాటసారి} mean? (\textquotedblleft{}A wayfarer\textquotedblright{}.) Say it once more.
\end{tcolorbox}
